\subsection{What is a logic?}
\label{sec:what-is-a-logic}
\slidesrc{Propositional logic: 1--3}

When most people say ``logic'' they mean either \emph{propositional logic} or
\emph{first-order predicate logic}. In fact the word is much broader: philosophers,
computer scientists and mathematicians have studied literally hundreds of logics (temporal
logics, modal logics, logics of belief, \ldots). What do they have in common?

\begin{nbdef}{Logic (slide 2)}
Any formal system can be considered a logic if it has:
\begin{itemize}[nosep]
  \item a well-defined \textbf{syntax}: it defines the syntactically acceptable objects of
        the language, called \textbf{well-formed formulae} (wff), or simply
        \textbf{formulae};
  \item a well-defined \textbf{semantics}: it associates each formula with a
        \textbf{meaning};
  \item a well-defined \textbf{proof theory}: it is concerned with \textbf{manipulating}
        formulae according to certain rules.
\end{itemize}
\end{nbdef}

\begin{center}
\begin{tikzpicture}[font=\small, >=Stealth,
  b/.style={draw, rounded corners, thick, align=center, text width=4.1cm, minimum height=1.8cm}]
  \node[b, draw=nbblue, fill=nbblue!6] (syn) at (0,0)
    {\textbf{Syntax}\\[2pt] which strings are formulas?\\[2pt]\footnotesize
     $p \land (q \limp r)$ \textcolor{nbgreen}{yes}\quad $p \land \limp$ \textcolor{nbred}{no}};
  \node[b, draw=nbgreen, fill=nbgreen!6] (sem) at (5,0)
    {\textbf{Semantics}\\[2pt] what does a formula mean? when is it true?\\[2pt]\footnotesize truth tables, valuations};
  \node[b, draw=nbpurple, fill=nbpurple!6] (pro) at (10,0)
    {\textbf{Proof theory}\\[2pt] how do we derive formulas by rules?\\[2pt]\footnotesize natural deduction};
  \draw[->, thick] (syn) -- (sem);
  \draw[->, thick] (syn.south) to[bend right=18] (pro.south);
\end{tikzpicture}
\captionof{figure}{The three ingredients of every logic. Semantics and proof theory are two
independent ways of talking about formulas: one with \emph{truth}, one with \emph{rules}.
A good proof theory derives exactly the formulas that are semantically valid.}
\end{center}

\begin{nbintuition}[A programming-language analogy]
Think of a programming language. Its \emph{syntax} says which programs compile
(\code{x = 3 +;} does not). Its \emph{semantics} says what a compiling program does when it
runs. A \emph{proof theory} is like a set of safe refactoring rules: you can transform the
code step by step without ever running it, and the rules guarantee that the meaning is
preserved. Logic works the same way, with formulas instead of programs and truth instead
of behaviour.
\end{nbintuition}

\begin{nbwhy}
Every symbolic AI system makes these three choices. A Prolog program (Part~IV) has a syntax
(facts and rules), a semantics (what is logically implied by the program) and a proof
procedure (SLD resolution) that the interpreter runs. A BDI agent (Part~V) stores its
beliefs as logical formulas and uses the same machinery to decide which plans apply.
\end{nbwhy}
